\documentclass[12pt,a4paper]{article}

\usepackage[utf8]{inputenc}
\usepackage[T1]{fontenc}
\usepackage{amsmath,amssymb,amsthm,mathtools}
\usepackage[colorlinks=true,linkcolor=blue,citecolor=red,urlcolor=blue]{hyperref}
\usepackage{geometry}
\geometry{margin=2.5cm}
\usepackage{enumitem}

\title{The FST-RH Framework: A Gauge-theoretic Reformulation of the Riemann Hypothesis\\ \large Executive Summary and Consolidated Thesis\thanks{Zenodo project DOI: 10.5281/zenodo.RH\_DOI\_PLACEHOLDER}}
\author{Lukas Geiger\\
\small Bernau, Germany\\
\small ORCID: \href{https://orcid.org/0009-0005-7296-1534}{0009-0005-7296-1534}}
\date{March 15, 2026}

\begin{document}
\maketitle

\noindent\textbf{MSC 2020:} 11M26, 11M36, 82B05, 60F05\\
\textbf{Keywords:} Riemann Hypothesis, gauge theory, transfer operators, Li coefficients, spectral certificates

\section{Overview: The FST-RH Trilogy}

The FST-RH framework develops a gauge-theoretic perspective on the Riemann Hypothesis (RH), connecting thermodynamic formalism on primes to the Li coefficients via the completed zeta function $\xi(s)$. This trilogy explores the bridge between analytic number theory and operator-theoretic stability.

\begin{description}
    \item[Part I: Prime Hub Graphs and the Vacuum Trace] constructs the motivational bridge. It defines the \emph{Prime Hub Graph} transfer operator and argues that the zeros of $\zeta(s)$ appear as its resonances. It introduces the Maximum Entropy Production Principle (MEPP) as a heuristic driver for critical-line symmetry.
    \item[Part II: Analytic Obstructions and the Limits of Monotonicity] establishes the Bombieri--Lagarias decomposition $\lambda_n = \lambda_n^{(\infty)} + \lambda_n^{(\mathrm{fin})}$ and identifies the ``Archimedean Dominance Paradox'': the archimedean Gamma-block provides global convexity, while the arithmetic prime-sum part $F(\sigma)$ changes sign at $\sigma_0 \approx 1.14$. It shows that Gibbs-normalized prime-sum formulations capture only local structure near $\sigma = 1$ and cannot serve as a global proxy for the Li coefficients.
    \item[Part III: Spectral Analysis and Certificate Structure] analyzes the spectral structure of the symmetrized arithmetic kernel, finding a stable negative inertia index of $2$. It develops the S-procedure certificate framework and discusses the scope and limitations of the gauge-correction approach.
\end{description}

\section{Key Results and Honest Assessment}

\subsection{What is Rigorously Established}
\begin{enumerate}
    \item \textbf{Bombieri--Lagarias Decomposition} (Part II): $\lambda_n = \lambda_n^{(\infty)} + \lambda_n^{(\mathrm{fin})}$ via the completed xi function $\xi(s)$, with $\lambda_n^{(\mathrm{fin})}$ computed from $g(s) = (s-1)\zeta(s)$, which is holomorphic at $s=1$ \cite{BombieriLagarias1999}.
    \item \textbf{Convergence of Prime-Sum Functionals} (Part II): $A(\sigma), C(\sigma), D(\sigma)$ converge absolutely at $\sigma = 1$. The target functional $F(\sigma) = A B - P(C+D)$ is positive for $\sigma \in (1, 1.14)$ and negative for $\sigma > 1.14$.
    \item \textbf{Spectral Census} (Part III): The symmetrized arithmetic kernel $\Ksym$ has a stable negative inertia index of $2$ for $N \ge 300$ primes and $\sigma \in (1, 5)$.
    \item \textbf{Gauge Equivalence} (Part III): For each fixed $\sigma$, the S-procedure provides a tautological equivalence between $F(\sigma) \ge 0$ and the existence of an admissible PSD gauge.
\end{enumerate}

\subsection{What Does NOT Work}
\begin{enumerate}
    \item \textbf{Global Gauge Reduction}: The claim ``RH $\Leftrightarrow$ $F(\sigma) \ge 0$ for all $\sigma > 1$'' is false, since $F(\sigma) < 0$ for $\sigma > 1.14$ independently of RH. The Gibbs-normalized prime-sum formulation cannot serve as a global proxy for the Li coefficients.
    \item \textbf{Gibbs Normalization}: The expectation $\mathbb{E}_{\mu_\sigma}[X_{n,\sigma}] = -A(\sigma)/P(\sigma)$ annihilates the finite-part contribution as $\sigma \to 1^+$ (the $1/P(\sigma)$ normalization kills the finite rest). The ``Critical Identification'' in the Gibbs framework converges to $0$, not $\lambda_n$.
    \item \textbf{Monotonicity Argument}: The true function $\frac{d}{ds}\log\xi(\sigma)$ is \emph{increasing} (convex, $\frac{d^2}{ds^2}\log\xi > 0$), not decreasing. The FST observable $m_n(\sigma)$ is decreasing but has the wrong limit.
\end{enumerate}

\subsection{Open Path Forward}
The core insight of the FST-RH program --- that the Li coefficients decompose into competing archimedean and arithmetic contributions, with archimedean dominance ensuring positivity --- remains valid as a structural observation. Numerical analysis reveals a three-regime structure:
\begin{itemize}
  \item \textbf{Small $n$ ($n \lesssim 8$):} $\lambda_n^{(\infty)} < 0$ but $\lambda_n^{(\mathrm{fin})} > 0$ compensates. Positivity is a delicate cancellation.
  \item \textbf{Transition ($n \approx 8$):} $\lambda_n^{(\infty)}$ crosses zero; both parts contribute positively.
  \item \textbf{Large $n$ ($n \gg 10$):} $\lambda_n^{(\infty)} \sim \frac{n}{2}\log n$ dominates; positivity is automatic.
\end{itemize}
Under RH, $|1-1/\rho| = 1$ for all zeros, giving $\lambda_n = \sum_\rho [1 - \cos(n\theta_\rho)] \ge 0$. Without RH, off-line zeros produce exponentially growing negative terms. Proving $\lambda_n \ge 0$ for all $n$ is therefore equivalent to RH \cite{Li1997}. Five concrete open problems are formulated in Part~IV.

\subsection{Connection to Connes' Program}
Part~IV additionally provides numerical evidence for the open problem identified in Connes' 2025 survey \cite{Connes2025}: that the smallest eigenvalue of the Weil quadratic form $\mathrm{QW}_\lambda$ is simple with even eigenfunction (Theorem~6.1). Using the corrected archimedean kernel $e^{u/2}/(2\sinh u)$ and Cauchy interlacing in a cosine Galerkin basis, we verify that the even-sector spectral gap is strictly positive for all $\lambda \in [5,200]$. The even--odd sector comparison confirms even dominance for $\lambda \ge 30$, with the crossover from odd to even ground state at $\lambda \approx 28$.

\section{Final Status}
The FST-RH program provides a thermodynamic and gauge-theoretic perspective on the Riemann Hypothesis, identifying several structural features (archimedean dominance, spectral census, boundary-layer positivity) and producing the strongest known numerical evidence for Connes' Theorem~6.1 prerequisites. It does \emph{not} constitute a proof or a reduction to a simpler problem. The framework's value lies in its conceptual organization, the identification of precise analytical obstructions, and the numerical verification of spectral conditions for the Weil quadratic form.

\bibliographystyle{plain}
\bibliography{fst-rh-references}

\end{document}
